\documentclass[12pt,a4paper]{article}

% ---------- Langue et encodage ----------
\usepackage[french]{babel}
%\usepackage[T1]{fontenc}
%\usepackage[utf8]{inputenc}

% ---------- Mise en page ----------
\usepackage[margin=2.2cm]{geometry}
\usepackage{array,booktabs,tabularx}
\usepackage{amsmath,amssymb}
\usepackage{xcolor}
\usepackage{tcolorbox}
\tcbuselibrary{skins,breakable}
\usepackage{enumitem}
\usepackage{graphicx}
\usepackage{tikz}
\usepackage{pgfplots}
\pgfplotsset{compat=1.18}
\usepackage{fancyhdr}
\usepackage{makecell}
\usepackage{pgffor}
\usepackage{currfile}
\usepackage{lastpage}

\graphicspath{ {../../Images/} }
\input{\currfiledir ../def.tex}
\def \Document {Activité 1}

% ---------- Couleurs ----------
\definecolor{bleuF}{RGB}{47,87,163}   % Femmes
\definecolor{orH}{RGB}{242,170,60}    % Hommes
\definecolor{cadre}{RGB}{47,87,163}

% ---------- Boîtes ----------
\newtcolorbox{objectif}{
  colback=bleuF!6, colframe=cadre, boxrule=0.6pt, arc=2pt,
  left=6pt, right=6pt, top=4pt, bottom=4pt,
  title={\textbf{Objectif}}, fonttitle=\bfseries,
  coltitle=cadre, colbacktitle=bleuF!12
}
\newtcolorbox{bilan}{
  colback=orH!8, colframe=cadre, boxrule=0.6pt, arc=2pt,
  left=6pt, right=6pt, top=4pt, bottom=4pt,
  title={\textbf{Bilan}}, fonttitle=\bfseries,
  coltitle=cadre, colbacktitle=orH!20,
  breakable
}
\newtcolorbox{correctionbox}{
  colback=green!5, colframe=green!40!black, boxrule=0.6pt, arc=2pt,
  left=6pt, right=6pt, top=4pt, bottom=4pt,
  breakable
}

% \checkmark symbol fallback (amssymb provides \checkmark)

% ---------- Ligne de réponse pointillée ----------
\newcommand{\reponse}[1]{%
  \par\vspace{2mm}
  \foreach \i in {1,...,#1}{\rule{0pt}{0pt}\dotfill\\[3mm]}%
}

\begin{document}


\pagestyle{fancy}

\fancyhf{} % clear existing header/footer entries
\fancyhead[R]{ {\small \Niveau}}
\fancyhead[L]{ {\small \ChapitreCourt}}
\fancyhead[C]{ {\small \Document}}
%\fancyhead[R]{\large \textbf{\textbf{\Document\ (B) -- 1h}}}
\fancyfoot[L]{ A. Fernandez }
\fancyfoot[R]{{\small \thepage\  / \pageref{LastPage}}}

\begin{objectif}
Compléter un tableau à double entrée à partir de proportions en pourcentages.
\end{objectif}

\bigskip
\noindent On souhaite construire un tableau représentant la répartition des ouvriers en France en 2019 par catégorie socioprofessionnelle et par sexe à partir des informations suivantes, données en pourcentage.

\bigskip
\noindent\textbf{Document 1 --- Répartition par sexe au sein de chaque catégorie}

\begin{center}
\begin{tikzpicture}
\begin{axis}[
  xbar stacked,
  width=0.74\textwidth, height=7.2cm,
  xmin=0, xmax=100,
  xlabel={\%},
  bar width=6.5mm,
  %enlarge y limits=0.11,
  symbolic y coords={Agricoles, Non qual.\ artisanal, Non qual.\ industriel, Manut.\ / magasinage / transport, Chauffeurs, Qual.\ artisanal, Qual.\ industriel},
  ytick=data,
  y tick label style={font=\footnotesize, align=right},
  axis y line*=none, axis x line*=bottom, tickwidth=0pt,
  nodes near coords, nodes near coords align={center},
  every node near coord/.append style={font=\scriptsize, white},
  legend style={at={(0.5,-0.20)}, anchor=north, legend columns=-1, draw=none, font=\small},
  ]
\addplot[fill=bleuF] coordinates {
  (10,Qual.\ industriel) (20,Qual.\ artisanal) (10,Chauffeurs)
  (20,Manut.\ / magasinage / transport) (20,Non qual.\ industriel)
  (30,Non qual.\ artisanal) (40,Agricoles)};
\addplot[fill=orH] coordinates {
  (90,Qual.\ industriel) (80,Qual.\ artisanal) (90,Chauffeurs)
  (80,Manut.\ / magasinage / transport) (80,Non qual.\ industriel)
  (70,Non qual.\ artisanal) (60,Agricoles)};
\legend{Femmes, Hommes}
\end{axis}
\end{tikzpicture}

%{\footnotesize Source~: Insee.}
\end{center}

\bigskip
\noindent\textbf{Document 2 --- Répartition des ouvriers par catégorie socioprofessionnelle}

\begin{center}
\begin{tabular}{lc}
\toprule
\textbf{Catégorie socioprofessionnelle} & \textbf{\%} \\
\midrule
Ouvriers qualifiés de type industriel & 20 \\
Ouvriers qualifiés de type artisanal & 25 \\
Chauffeurs & 10 \\
Ouvriers qualifiés de la manutention, du magasinage et du transport & 5 \\
Ouvriers non qualifiés de type industriel & 15 \\
Ouvriers non qualifiés de type artisanal & 20 \\
Ouvriers agricoles & 5 \\
%Ouvriers qualifiés de type industriel & 19,9 \\
%Ouvriers qualifiés de type artisanal & 24,6 \\
%Chauffeurs & 12,5 \\
%Ouvriers qualifiés de la manutention, du magasinage et du transport & 8,4 \\
%Ouvriers non qualifiés de type industriel & 16,4 \\
%Ouvriers non qualifiés de type artisanal & 13,6 \\
%Ouvriers agricoles & 4,6 \\
\midrule
\textbf{Ensemble des ouvriers} & \textbf{100} \\
\bottomrule
\end{tabular}

%{\footnotesize Source~: Insee.}
\end{center}

\newpage
\noindent\textbf{Questions}

\begin{enumerate}[label=\textbf{\arabic*.}, leftmargin=1.3em]
\item Laquelle des affirmations suivantes est juste~? Cocher la bonne réponse.
\begin{itemize}[label=$\square$]
  \item « 40~\% des ouvriers sont des femmes ouvrières agricoles. »
  \item « 40~\% des ouvriers agricoles sont des femmes. »
\end{itemize}
Expliquer la différence entre ces deux affirmations.
\reponse{3}

\item Interpréter le nombre 10 qui apparaît dans la troisième ligne du second tableau.
\reponse{3}

\item À partir de ces tableaux, déterminer la proportion de femmes chauffeurs parmi l'ensemble des ouvriers, en pourcentage.
\reponse{3}
\end{enumerate}

\bigskip
\begin{bilan}
À l'aide des tableaux exposés dans cette activité, compléter le tableau ci-dessous avec des pourcentages calculés par rapport à l'ensemble des ouvriers de France. On arrondira au besoin les résultats à 0,1~\%.

\bigskip
\begin{center}
\renewcommand{\arraystretch}{3}
\resizebox{\textwidth}{!}{%
\begin{tabular}{|l|c|c|c|c|c|c|c|c|}
\hline
 & \makecell{Qual.\\industriel} & \makecell{Qual.\\artisanal} & Chauffeurs & \makecell{Manut. /\\magasinage /\\transport} & \makecell{Non qual.\\industriel} & \makecell{Non qual.\\artisanal} & Agricoles & \textbf{Ensemble} \\
\hline
\textbf{Femmes} & & & & & & & & \\
\hline
\textbf{Hommes} & & & & & & & & \\
\hline
\textbf{Ensemble} & & & & & & & & \\
\hline
\end{tabular}}
\end{center}
\end{bilan}


% \clearpage
% \pagestyle{empty}
% \begin{center}
% {\LARGE\textbf{Correction}}\\[2pt]
% {\large Tableau à double entrée : fréquences marginales et conditionnelles}
% \end{center}
% \vspace{3mm}
%  
% \begin{enumerate}[label=\textbf{\arabic*.}, leftmargin=1.3em]
%  
% \item \textbf{Réponse~:} c'est la \emph{seconde} affirmation qui est juste~: « 40~\% des ouvriers agricoles sont des femmes. »
%  
% \begin{correctionbox}
% \textbf{Explication.} Le nombre 40~\% est une \emph{fréquence conditionnelle}~: c'est le pourcentage de femmes calculé \emph{parmi les seuls ouvriers agricoles} (la catégorie « ouvriers agricoles » sert de référence, de « dénominateur »).
%  
% La première affirmation confond cette fréquence conditionnelle avec une \emph{fréquence par rapport à l'ensemble des ouvriers} (une fréquence « croisée », ou jointe)~: le pourcentage de femmes ouvrières agricoles parmi \emph{tous} les ouvriers n'est pas 40~\%, mais se calcule à la question~3 (on trouve 2~\%, bien plus faible), car les ouvriers agricoles ne représentent eux-mêmes que 5~\% de l'ensemble des ouvriers.
% \end{correctionbox}
%  
% \item \textbf{Réponse~:} 10 est une \emph{fréquence marginale}~: parmi l'ensemble des ouvriers de France (hommes et femmes confondus), 10~\% exercent la profession de chauffeur.
%  
% \item \textbf{Réponse~:} On cherche la fréquence \emph{croisée} (femmes $\cap$ chauffeurs) par rapport à l'ensemble des ouvriers. Elle s'obtient en multipliant la fréquence marginale de la catégorie « chauffeurs » par la fréquence conditionnelle de femmes \emph{parmi} les chauffeurs~: c'est 10~\% de 10, c'est-à-dire
% \[
% 10\,\% \times 10 = \boxed{1\,\%}
% \]
% \end{enumerate}
%  
% \bigskip
% \begin{correctionbox}
% \textbf{Méthode générale (Bilan).} Pour chaque catégorie, la fréquence croisée par rapport à l'ensemble des ouvriers s'obtient par~:
% \begin{center}
% \itshape fréquence croisée $=$ fréquence marginale $\times$ fréquence conditionnelle (Femmes ou Hommes)
% \end{center}
% Les fréquences croisées « Hommes » s'en déduisent aussi par différence à la fréquence marginale de la catégorie (ou en multipliant par la fréquence conditionnelle complémentaire).
% \end{correctionbox}
%  
% \bigskip
% \begin{center}
% \resizebox{\textwidth}{!}{%
% \begin{tabular}{|l|c|c|c|c|c|c|c|c|}
% \hline
%  & \makecell{Qual.\\industriel} & \makecell{Qual.\\artisanal} & Chauffeurs & \makecell{Manut. /\\magasinage /\\transport} & \makecell{Non qual.\\industriel} & \makecell{Non qual.\\artisanal} & Agricoles & \textbf{Ensemble} \\
% \hline
% \textbf{Femmes} & 2 & 5 & 1 & 1 & 3 & 6 & 2 & \textbf{20} \\
% \hline
% \textbf{Hommes} & 18 & 20 & 9 & 4 & 12 & 14 & 3 & \textbf{80} \\
% \hline
% \textbf{Ensemble} & \textbf{20} & \textbf{25} & \textbf{10} & \textbf{5} & \textbf{15} & \textbf{20} & \textbf{5} & \textbf{100} \\
% \hline
% \end{tabular}}
% \end{center}
%  
% \bigskip
% {\footnotesize\textit{Remarque~:} la dernière ligne (« Ensemble ») redonne exactement les fréquences marginales du Document~2, et la dernière colonne redonne la répartition marginale par sexe (20~\% de femmes, 80~\% d'hommes parmi l'ensemble des ouvriers), ce qui permet de vérifier la cohérence du tableau.}

\end{document}